\documentclass[11pt]{article}
% ===========================================================================
%  Shared preamble for the AP Calculus AB exam-preparation worksheets.
%  Input right after \documentclass. Files beginning with "_" are not compiled
%  on their own. Navy and gold, to match the course's lesson pages.
% ===========================================================================
\usepackage[letterpaper,top=2.3cm,bottom=2.2cm,left=1.7cm,right=1.7cm,headheight=28pt,footskip=30pt]{geometry}
\usepackage{amsmath,amssymb}
\usepackage[T1]{fontenc}
\usepackage[utf8]{inputenc}
\usepackage{lmodern}
\usepackage{xcolor}
\usepackage[most]{tcolorbox}
\usepackage{enumitem}
\usepackage{fancyhdr}
\usepackage{array}
\usepackage{tikz}
\usepackage{pgfplots}
\pgfplotsset{compat=1.18}

\definecolor{navy}{HTML}{161E54}
\definecolor{ink}{HTML}{27357E}
\definecolor{gold}{HTML}{B98A00}
\definecolor{goldbg}{HTML}{FFFDF3}
\definecolor{inkbg}{HTML}{F6F7FD}
\definecolor{rule}{HTML}{C9CFF0}
\definecolor{okgreen}{HTML}{1B7A44}
\definecolor{okbg}{HTML}{F1F8F3}

\pagestyle{fancy}
\fancyhf{}
\fancyhead[L]{\footnotesize Name: \underline{\hspace{4cm}}}
\fancyhead[C]{\textbf{\color{navy}Integration Academy}\\[-2pt]{\scriptsize Calculus AB \textperiodcentered{} AP\textsuperscript{\textregistered} Exam Preparation}}
\fancyfoot[L]{\tiny\color{black!55}\parbox[t]{0.78\textwidth}{AP\textsuperscript{\textregistered} is a registered trademark of the College Board, which was not involved in the production of, and does not endorse, this material. \textcopyright{} Integration Academy.}}
\fancyfoot[R]{\scriptsize\thepage}
\renewcommand{\headrulewidth}{1.2pt}
\renewcommand{\headrule}{\hbox to\headwidth{\color{gold}\leaders\hrule height \headrulewidth\hfill}}
\renewcommand{\footrulewidth}{0pt}
\setlength{\parindent}{0pt}
\setlength{\parskip}{4pt}

% Title block: \wstitle{code}{title}{unit line}
\newcommand{\wstitle}[3]{\begin{center}{\small\bfseries\color{gold}WORKSHEET #1}\\[3pt]{\LARGE\bfseries\color{navy}#2}\\[3pt]{\small\color{black!60}#3}\end{center}\vspace{2pt}}
% Section heading
\newcommand{\sech}[1]{\par\vspace{10pt}{\large\bfseries\color{navy}#1}\par\nobreak\vspace{2pt}{\color{rule}\hrule height0.8pt}\vspace{7pt}}
% A level tag printed beside a group of questions
\newcommand{\level}[1]{\par\vspace{6pt}{\footnotesize\bfseries\color{gold}\MakeUppercase{#1}}\par\nobreak\vspace{2pt}}

\newtcolorbox{ideabox}{colframe=ink,colback=inkbg,coltitle=white,colbacktitle=ink,fonttitle=\bfseries\small,title={Key ideas},arc=2pt,boxrule=0.8pt,left=8pt,right=8pt}
\newtcolorbox{example}[1]{colframe=rule,colback=white,coltitle=navy,colbacktitle=inkbg,fonttitle=\bfseries\small,title={#1},arc=2pt,boxrule=0.8pt,borderline west={3pt}{0pt}{ink},left=10pt,right=8pt}
\newtcolorbox{frbox}[1]{colframe=gold,colback=goldbg,coltitle=white,colbacktitle=gold,fonttitle=\bfseries\small,title={#1},arc=2pt,boxrule=0.8pt,breakable,left=8pt,right=8pt}
\newtcolorbox{solbox}{colframe=okgreen,colback=okbg,coltitle=white,colbacktitle=okgreen,fonttitle=\bfseries\small,title={Worked solutions},arc=2pt,boxrule=0.8pt,breakable,left=8pt,right=8pt}
\newtcolorbox{scorebox}[1]{colframe=navy,colback=white,coltitle=white,colbacktitle=navy,fonttitle=\bfseries\small,title={#1},arc=2pt,boxrule=0.8pt,breakable,left=8pt,right=8pt}

\newcommand{\soln}{\par\smallskip\textbf{\color{ink}Solution.}\ }
\newcommand{\result}[1]{\par\smallskip\textbf{#1}}
% blank working space of a given height
\newcommand{\work}[1]{\par\vspace{#1}}
% a point value in the margin of a scoring guide
\newcommand{\pts}[1]{\hfill{\footnotesize\bfseries\color{gold}[#1]}}

% Coordinate grid for graph questions: \begin{apgraph}{xmin}{xmax}{ymin}{ymax} ... \end{apgraph}
\newenvironment{apgraph}[4]{\begin{center}\begin{tikzpicture}\begin{axis}[
  axis lines=middle, xmin=#1, xmax=#2, ymin=#3, ymax=#4,
  grid=both, grid style={black!12}, xtick distance=1, ytick distance=1,
  ticklabel style={font=\scriptsize}, xlabel={$x$}, ylabel={$y$},
  width=10cm, height=7.2cm, samples=80, clip=false,
  every axis plot/.append style={line width=1.1pt, color=ink}]}{\end{axis}\end{tikzpicture}\end{center}}
\newcommand{\closeddot}[2]{\addplot[only marks, mark=*, mark size=2.2pt, color=ink] coordinates {(#1,#2)};}
\newcommand{\opendot}[2]{\addplot[only marks, mark=*, mark size=2.2pt, color=ink, mark options={fill=white}] coordinates {(#1,#2)};}

% Multiple-choice options: four across, or stacked when the options are long.
\newcommand{\opts}[4]{\par\vspace{3pt}\begin{tabular}{@{}p{0.235\linewidth}p{0.235\linewidth}p{0.235\linewidth}p{0.235\linewidth}@{}}(A)\ #1 & (B)\ #2 & (C)\ #3 & (D)\ #4\end{tabular}\par\vspace{9pt}}
\newcommand{\optsv}[4]{\par\vspace{3pt}\begin{tabular}{@{}l@{\ \ }p{0.9\linewidth}@{}}(A) & #1\\(B) & #2\\(C) & #3\\(D) & #4\end{tabular}\par\vspace{9pt}}

\fancyhead[R]{\footnotesize\color{black!60}Worksheet 8.10}
\begin{document}

\wstitle{8.10}{Unit 8 Synthesis: The Area and Volume Free-Response Question}{Unit 8: Applications of Integration}

\begin{ideabox}
\begin{itemize}[leftmargin=*,itemsep=2pt,topsep=2pt]
  \item The usual parts: an area, a volume of revolution, and a volume with known cross-sections, all for one region.
  \item Find the points of intersection first, and store them in the calculator. Do not round until the final answer.
  \item Write every integral in full, with limits, before giving its value. The set-up earns most of the points.
  \item \textquotedblleft{}Write, but do not evaluate\textquotedblright{} means stop at the integral.
\end{itemize}
\end{ideabox}

\sech{Worked examples}
\begin{example}{Example 1 \textperiodcentered{} The area}
Let $R$ be the region enclosed by $y=\sqrt{x}$ and $y=\dfrac x3$. Find the area of $R$.\soln They meet at $x=0$ and $x=9$: $\displaystyle\int_{0}^{9}\left(\sqrt x-\tfrac x3\right)\,dx=18-\dfrac{27}{2}$.\result{$\dfrac92$.}
\end{example}

\begin{example}{Example 2 \textperiodcentered{} A solid of revolution}
Find the volume when $R$ is revolved about the $x$-axis.\soln $\pi\displaystyle\int_{0}^{9}\left(x-\tfrac{x^2}{9}\right)\,dx=\pi\left(\dfrac{81}{2}-27\right)$.\result{$\dfrac{27\pi}{2}$.}
\end{example}

\begin{example}{Example 3 \textperiodcentered{} Known cross-sections}
$R$ is the base of a solid whose cross-sections perpendicular to the $x$-axis are squares. Find the volume.\soln $\displaystyle\int_{0}^{9}\left(\sqrt x-\tfrac x3\right)^2\,dx=\dfrac{81}{2}-\dfrac{324}{5}+27$.\result{$\dfrac{27}{10}$.}
\end{example}

\newpage
\sech{Practice}
\emph{A calculator may be used where a decimal answer is requested. Give decimals to three places.}
\level{Core}
\begin{enumerate}[leftmargin=*,itemsep=2pt]
  \item Let $S$ be the region enclosed by $y=x^2$ and $y=3x$. Find the area of $S$.\work{2.3cm}
  \item Write, but do not evaluate, an integral for the volume when $S$ is revolved about the $x$-axis.\work{2.3cm}
  \item Evaluate the integral in Question 2.\work{2.3cm}
  \item $S$ is the base of a solid whose cross-sections perpendicular to the $x$-axis are squares. Write, but do not evaluate, an integral for the volume.\work{2.3cm}
\end{enumerate}

\level{Exam level}
\begin{enumerate}[leftmargin=*,itemsep=2pt,start=5]
  \item Let $T$ be the region enclosed by $y=4-x^2$ and $y=e^{x}$. Find the $x$-coordinates of the two points of intersection.\work{2.8cm}
  \item Find the area of $T$.\work{2.8cm}
  \item Find the volume of the solid generated when $T$ is revolved about the $x$-axis.\work{2.8cm}
  \item $T$ is the base of a solid whose cross-sections perpendicular to the $x$-axis are semicircles with diameter in the base. Find the volume.\work{2.8cm}
\end{enumerate}

\begin{frbox}{Question 9 \textperiodcentered{} Free-response practice}
Let $R$ be the region enclosed by $y=2\sqrt{x}$ and $y=x$.
\begin{enumerate}[label=(\alph*),leftmargin=*,itemsep=2pt]
  \item Find the area of $R$.\work{2.6cm}
  \item Find the volume of the solid generated when $R$ is revolved about the $x$-axis.\work{2.6cm}
  \item $R$ is the base of a solid whose cross-sections perpendicular to the $y$-axis are semicircles with diameter in the base. Find the volume.\work{2.6cm}
\end{enumerate}
\end{frbox}

\newpage
\begin{solbox}
\begin{enumerate}[leftmargin=*,itemsep=4pt]
  \item $\displaystyle\int_{0}^{3}(3x-x^2)\,dx=\dfrac92$.
  \item $\pi\displaystyle\int_{0}^{3}(9x^2-x^4)\,dx$.
  \item $\pi\left(81-\dfrac{243}{5}\right)=\dfrac{162\pi}{5}$.
  \item $\displaystyle\int_{0}^{3}(3x-x^2)^2\,dx$.
  \item $x\approx-1.965$ and $x\approx1.058$.
  \item $\displaystyle\int_{-1.965}^{1.058}\big(4-x^2-e^{x}\big)\,dx\approx6.428$.
  \item $\pi\displaystyle\int_{-1.965}^{1.058}\big[(4-x^2)^2-e^{2x}\big]\,dx\approx84.705$.
  \item $\dfrac\pi8\displaystyle\int_{-1.965}^{1.058}\big(4-x^2-e^{x}\big)^2\,dx\approx6.422$.
  \item \begin{enumerate}[label=(\alph*),leftmargin=*,itemsep=3pt]
    \item They meet at $x=0$ and $x=4$: $\displaystyle\int_{0}^{4}(2\sqrt x-x)\,dx=\dfrac{32}{3}-8=\dfrac83$.
    \item $\pi\displaystyle\int_{0}^{4}(4x-x^2)\,dx=\pi\left(32-\dfrac{64}{3}\right)=\dfrac{32\pi}{3}$.
    \item At height $y$ the region runs from $x=\dfrac{y^2}{4}$ to $x=y$: $\dfrac\pi8\displaystyle\int_{0}^{4}\left(y-\tfrac{y^2}{4}\right)^2\,dy=\dfrac\pi8\cdot\dfrac{32}{15}=\dfrac{4\pi}{15}$.
  \end{enumerate}
\end{enumerate}
\end{solbox}

\end{document}
